\section{The absorption models in full}
\label{m:app:absorption}

\Cref{m:sec:moc} developed the method of characteristics on one absorption model.
This appendix collects the family that model belongs to: the simplest member, in
which particles are only taken up; and the hardest, in which they may replicate
once inside, and for which the closed form obtained is a result of this thesis.

The setting is a single compartment drawing particles from a surrounding medium.
Once several compartments share a medium, a question arises that the models of
\cref{m:sec:chains} do not address: what becomes of the contents of the first
compartment to rupture? The shared medium acquires, at a stroke, a potentially
large cohort of particles available for subsequent absorption by the compartments
that remain. The analysis elsewhere in this thesis assumes \emph{immediate
transfer}, in which the entire contents of a ruptured compartment relocate at once.
The alternative is an \emph{absorption period}, during which absorption events
redistribute the released cohort among the remaining compartments and only then do
birth, death and rupture resume.

Throughout, $X_t$ and $Y_t$ count particles inside and outside the compartment, the
initial condition is $X_0=0$, $Y_0=N$, and events have exponentially distributed
inter-event times.

\subsection{Absorption only}
\label{m:app:absonly}

The simplest case has no dynamics inside the compartment at all:
\begin{align*}
\pr{\Delta X=1,\ \Delta Y=-1\mid Y_t=j} &= j\alpha\,\Delta t, &&\text{(absorption)}\\
\pr{\Delta X=0,\ \Delta Y=0\mid Y_t=j} &= 1-j\alpha\,\Delta t, &&\text{(no event)}
\end{align*}
the factor $j$ recording that each exterior particle is independently subject to
the same probability $\alpha\Delta t$ of being taken up. The forward equations are
\begin{equation}
\label{m:eq:FK1}
\ddt{p_{i,j}} = \alpha(j+1)p_{i-1,j+1} - \alpha j\,p_{i,j},
\end{equation}
and the construction of \cref{m:sec:backward} gives
\begin{equation}
\label{m:eq:PDE1}
\pddt{G} = \alpha(x-y)\pdd{G}{y},
\qquad
G(x,y,0) = y^N .
\end{equation}
The sums involved are finite, because $p_{i,j}(t)=0$ whenever $i+j\neq N$: particles
in this model are neither created nor destroyed.

Because the particles never interact, the generating function factorises. Writing
$G_n$ for the function with $n$ initial exterior particles and decomposing
$X_t = \sum_kX_t^{(k)}$, $Y_t = \sum_kY_t^{(k)}$ into single-particle contributions,
independence gives
\begin{equation}
\label{m:eq:indRelPGF}
G_n(x,y,t)
= \ex{\prod_{k=1}^{n}x^{X_t^{(k)}}y^{Y_t^{(k)}}}
= \prod_{k=1}^{n}\ex{x^{X_t^{(k)}}y^{Y_t^{(k)}}}
= \bigl(G_1(x,y,t)\bigr)^n .
\end{equation}
For a single particle only two states are reachable, so
$\dd p_{0,1}/\dd t = -\alpha p_{0,1}$ and $\dd p_{1,0}/\dd t = \alpha p_{0,1}$,
giving $p_{0,1}(t) = \ee^{-\alpha t}$ and $p_{1,0}(t) = 1-\ee^{-\alpha t}$
(\cref{m:fig:abs1}); hence
\begin{equation}
\label{m:eq:Gabsonly}
G_n(x,y,t) = \Bigl(\lrb{1-\ee^{-\alpha t}}x+\ee^{-\alpha t}y\Bigr)^n .
\end{equation}
This satisfies \cref{m:eq:PDE1} with the required initial condition, and the form is
recognisable: $X_t\sim\text{Binomial}(n,1-\ee^{-\alpha t})$, as it must be, since
the $n$ particles are absorbed independently and each has been taken up by time $t$
with probability $1-\ee^{-\alpha t}$.

\begin{figure}[H]
\centering
\includegraphics[width=0.58\textwidth]{abs1.pdf}
\caption{State probabilities for the single-particle absorption-only model. The
particle begins outside the compartment ($p_{0,1}$, decaying exponentially at rate
$\alpha$) and ends inside it ($p_{1,0}$). There is nowhere else for it to go, so
the two curves sum to one at every time.}
\label{m:fig:abs1}
\end{figure}

\subsection{Absorption--birth--death}
\label{m:app:absbirthdeath}

Introducing replication for particles inside the compartment, the factorisation
\cref{m:eq:indRelPGF} survives but no longer helps: from a single initial particle,
birth events make infinitely many states reachable, so $G_1$ is an infinite sum.
The partial differential equation must be solved directly.

Birth at per-capita rate $\lambda$ is added to the absorption--death model of
\cref{m:sec:mocworked}:
\begin{align*}
\pr{\Delta X=1,\ \Delta Y=-1\mid Y_t=j} &= j\alpha\,\Delta t, &&\text{(absorption)}\\
\pr{\Delta X=1,\ \Delta Y=0\mid X_t=i} &= i\lambda\,\Delta t, &&\text{(birth)}\\
\pr{\Delta X=-1,\ \Delta Y=0\mid X_t=i} &= i\mu\,\Delta t, &&\text{(death)}\\
\pr{\Delta X=0,\ \Delta Y=0\mid X_t=i,\ Y_t=j} &= 1-(j\alpha+i\mu+i\lambda)\Delta t, &&\text{(no event)}
\end{align*}
with forward equation
\begin{equation}
\label{m:eq:FK3}
\ddt{p_{i,j}} = \alpha(j+1)p_{i-1,j+1} + \mu(i+1)p_{i+1,j} + \lambda(i-1)p_{i-1,j}
- (\alpha j+\mu i+\lambda i)p_{i,j},
\end{equation}
and generating-function equation
\begin{equation}
\label{m:eq:PDE3}
\pddt{G} = \bigl(\mu-(\mu+\lambda)x+\lambda x^2\bigr)\pdd{G}{x}
+ \alpha(x-y)\pdd{G}{y},
\qquad
G(x,y,0)=y^N ,
\end{equation}
whose $x$-coefficient is exactly the quadratic met in \cref{m:eq:bdpde} for the
linear birth--death process. The characteristic system
\cref{m:eq:genericcharacteristics} is
\begin{align}
\dda{G}{\tau} &= 0, \label{m:eq:ch21}\\
\dda{t}{\tau} &= -1, \label{m:eq:ch22}\\
\dda{x}{\tau} &= \mu-(\mu+\lambda)x+\lambda x^2, \label{m:eq:ch23}\\
\dda{y}{\tau} &= \alpha(x-y), \label{m:eq:ch24}
\end{align}
with the initial conditions \cref{m:eq:ICpar}.

\paragraph{The equations for $G$ and $t$.}
As in \cref{m:sec:mocworked}, and for the same universal reasons,
\begin{equation}
\label{m:eq:GetaN2}
G(\sigma,\eta,\tau) = \eta^N,
\qquad
t = -\tau .
\end{equation}

\paragraph{The equation for $x$, \eqref{m:eq:ch23}.}
This is the Riccati equation anticipated in \cref{m:app:riccati}. Its right-hand
side factorises, and the roots are elementary rather than the surd expression that
mechanical use of the quadratic formula suggests:
\begin{equation}
\label{m:eq:xpm}
\mu-(\mu+\lambda)x+\lambda x^2 = \lambda(x-x_+)(x-x_-),
\qquad
x_+ = 1,\quad x_- = \frac{\mu}{\lambda} .
\end{equation}
Name the combination that appears throughout,
\begin{equation}
\label{m:eq:b1}
b_1 = \lambda(x_+-x_-) = \lambda-\mu ,
\end{equation}
the net interior growth rate. Separating variables and applying \cref{m:eq:ICpar}
gives the characteristic in the compact form
\begin{equation}
\label{m:eq:xratio}
\frac{x-x_+}{x-x_-} = \lrb{\frac{\sigma-x_+}{\sigma-x_-}}\ee^{b_1\tau},
\end{equation}
or, solved for $x$,
\begin{equation}
\label{m:eq:xabd}
x(\tau) = \frac{x_+(\sigma-x_-) - x_-(\sigma-x_+)\ee^{b_1\tau}}
{(\sigma-x_-) - (\sigma-x_+)\ee^{b_1\tau}} .
\end{equation}

\paragraph{The equation for $y$, \eqref{m:eq:ch24}.}
With an integrating factor this is
\begin{equation}
\label{m:eq:yint}
\dda{}{\tau}\lrb{\ee^{\alpha\tau}y} = \alpha\,\ee^{\alpha\tau}x(\tau),
\end{equation}
and the integral on the right is where the difficulty lies. Rather than substitute
\cref{m:eq:xabd} and split the resulting quotient, it is cleaner to expand
$x(\tau)$ first. Writing $x-x_- = (x-x_+)+(x_+-x_-)$ in \cref{m:eq:xratio} and
solving for $x$ gives
\begin{equation}
\label{m:eq:xgeom}
x(\tau) = x_+ + (x_+-x_-)\frac{\zeta(\tau)}{1-\zeta(\tau)},
\qquad
\zeta(\tau) = \lrb{\frac{\sigma-x_+}{\sigma-x_-}}\ee^{b_1\tau},
\end{equation}
so that for $|\zeta|<1$,
$x(\tau) = x_+ + (x_+-x_-)\sum_{n\ge1}\zeta(\tau)^n$. Every term is now a pure
exponential in $\tau$ and integrates immediately. Writing
$Z = (\sigma-x_+)/(\sigma-x_-)$, so that $\zeta(\tau)=Z\ee^{b_1\tau}$,
\begin{align*}
\alpha\int\ee^{\alpha\tau}x(\tau)\,\dd\tau
&= \alpha x_+\!\int\!\ee^{\alpha\tau}\dd\tau
 + \alpha(x_+-x_-)\sum_{n=1}^{\infty}Z^n\!\int\!\ee^{(\alpha+nb_1)\tau}\dd\tau\\
&= \ee^{\alpha\tau}\lrbs{x_+ + \alpha(x_+-x_-)\sum_{n=1}^{\infty}\frac{\zeta^n}{\alpha+nb_1}}
 + \text{const} .
\end{align*}
Setting $\delta = \alpha/b_1$ and using $\alpha(x_+-x_-)/b_1 = \alpha/\lambda$ by
\cref{m:eq:b1},
\begin{equation}
\label{m:eq:Psidef}
\Psi(\zeta) = x_+ + \frac{\alpha}{\lambda}\sum_{n=1}^{\infty}\frac{\zeta^{\,n}}{n+\delta},
\qquad
\delta = \frac{\alpha}{b_1} = \frac{\alpha}{\lambda-\mu} .
\end{equation}
Integrating \cref{m:eq:yint} gives
$\ee^{\alpha\tau}y = \ee^{\alpha\tau}\Psi(\zeta(\tau))+\text{const}$, and the
initial condition $y(\sigma,\eta,0)=\eta$, at which $\zeta(0)=Z$, fixes the constant
as $\eta-\Psi(Z)$, so
\begin{equation}
\label{m:eq:etasolved}
\eta = \ee^{\alpha\tau}\Bigl[y-\Psi\bigl(\zeta(\tau)\bigr)\Bigr] + \Psi(Z) .
\end{equation}
Finally, put $\tau=-t$ and write $W = (x-x_+)/(x-x_-)$, so that $\zeta(\tau)=W$ and
$Z = W\ee^{b_1t}$ by \cref{m:eq:xratio}. Substituting into \cref{m:eq:etasolved} and
applying \cref{m:eq:GetaN2} gives the solution.

\begin{proposition}[Generating function for absorption--birth--death]
\label{m:prop:Gabd}
With $x_+=1$, $x_-=\mu/\lambda$, $b_1=\lambda-\mu$, $\delta=\alpha/b_1$,
$W = (x-x_+)/(x-x_-)$ and $\Psi$ as in \cref{m:eq:Psidef},
\begin{equation}
\label{m:eq:Gabd}
\boxed{\;
G(x,y,t) = \Bigl\{\ee^{-\alpha t}\bigl[y-\Psi(W)\bigr] + \Psi\bigl(W\ee^{b_1t}\bigr)\Bigr\}^{N} . \;}
\end{equation}
\end{proposition}

Three checks confirm the result. At $t=0$ the two $\Psi$ terms cancel and $G=y^N$,
the required initial condition. At $x=y=1$ we have $W=0$ and $\Psi(0)=x_+=1$, so
$G=1$ and probability is conserved. And setting $x=1$ alone gives
$G(1,y,t) = (\ee^{-\alpha t}y+1-\ee^{-\alpha t})^N$, so
$Y_t\sim\text{Binomial}(N,\ee^{-\alpha t})$ --- exactly right, since the exterior
particles are absorbed independently at rate $\alpha$ and what happens inside the
compartment cannot affect them. Beyond these, \cref{m:eq:Gabd} has been verified
numerically against direct integration of \cref{m:eq:FK3}, agreeing to around
$10^{-13}$ across sub- and supercritical interior regimes and for $N=1,2,3$.

\subsection{A parameter-regular representation}
\label{m:app:abdregular}

The form \cref{m:eq:Gabd} is compact, but its derivation separated an antiderivative
into pieces that are individually singular at certain parameter values. The
underlying probability calculation supplies a second representation that stays
finite everywhere.

Let $F(x,t)$ be the generating function, at time $t$, of an interior linear
birth--death process started from a single particle. Conditioning on the first
event gives its backward equation,
\begin{equation}
\label{m:eq:interiorbackward}
\pddt{F} = \lambda F^2 - (\lambda+\mu)F + \mu,
\qquad F(x,0)=x ,
\end{equation}
which is \cref{m:eq:bdpde} read along characteristics. For $\lambda\neq\mu$ its
solution is the ratio already met in \cref{m:eq:xratio},
\begin{equation}
\label{m:eq:interiorratio}
\frac{F(x,t)-1}{F(x,t)-x_-} = \frac{x-1}{x-x_-}\,\ee^{b_1t},
\end{equation}
while at $\lambda=\mu$ the roots coincide and
\begin{equation}
\label{m:eq:interiorcritical}
F(x,t) = 1-\frac{1-x}{1+\lambda t(1-x)} .
\end{equation}
Now follow a single exterior particle. Either it has not been absorbed by time $t$,
which happens with probability $\ee^{-\alpha t}$ and contributes $y$; or it is
absorbed at some time $v\in(0,t)$, with density $\alpha\ee^{-\alpha v}$, and
thereafter founds an interior birth--death process of duration $t-v$. Conditioning
on the absorption time,
\begin{equation}
\label{m:eq:abdconvolution}
g(x,y,t) = \ee^{-\alpha t}y
+ \alpha\ee^{-\alpha t}\int_0^t\ee^{\alpha v}F(x,v)\,\dd v ,
\qquad
G(x,y,t) = g(x,y,t)^N ,
\end{equation}
the second identity because the $N$ exterior particles are independent throughout.
\Cref{m:eq:abdconvolution} is real and finite for $x,y\in[0,1]$ and all positive
rates, including $\lambda=\mu$ where \cref{m:eq:Gabd} degenerates. For
$\lambda\neq\mu$, substituting \cref{m:eq:interiorratio} into the integral and
integrating term by term returns \cref{m:eq:Gabd}; so \cref{m:eq:abdconvolution}
also specifies unambiguously which continuation of the special-function expression
is intended.

\subsection{Hypergeometric form, domain and resonance}
\label{m:app:abdhyp}

The series in \cref{m:eq:Psidef} is a Gauss hypergeometric function in disguise.
With
\begin{equation}
\label{m:eq:2F1def}
{}_2F_1(a,b;c;z) = \sum_{n=0}^{\infty}\frac{(a)_n(b)_n}{(c)_n}\frac{z^n}{n!},
\qquad
(q)_n = \begin{cases}1, & n=0,\\ q(q+1)\cdots(q+n-1), & n\ge1,\end{cases}
\end{equation}
where $(q)_n$ is the rising Pochhammer symbol \cite{Olver2010}, the case $a=1$,
$b=\delta$, $c=\delta+1$ collapses to a single ratio. Since $(1)_n=n!$ and
\begin{equation}
\label{m:eq:pochcancel}
\frac{(\delta)_n}{(\delta+1)_n}
= \frac{\delta(\delta+1)\cdots(\delta+n-1)}{(\delta+1)\cdots(\delta+n-1)(\delta+n)}
= \frac{\delta}{\delta+n},
\end{equation}
we have ${}_2F_1(1,\delta;\delta+1;z) = \sum_{n\ge0}\frac{\delta}{\delta+n}z^n$.
Comparing with \cref{m:eq:Psidef} and using
$\alpha/(\lambda\delta) = b_1/\lambda = x_+-x_-$,
\begin{equation}
\label{m:eq:Psihyp}
\Psi(\zeta) = x_+ + (x_+-x_-)\Bigl[{}_2F_1(1,\delta;\delta+1;\zeta)-1\Bigr] .
\end{equation}
\Cref{m:eq:Gabd} together with \cref{m:eq:Psihyp} is the closed form. It is not
pretty, but it is a genuine closed form in terms of a standard special function,
and the elementary series \cref{m:eq:Psidef} is what one would actually evaluate. An
alternative derivation, by way of the integral
$\int\ee^{h\tau}/(P+Q\ee^{k\tau})\,\dd\tau$, is given in \cref{m:app:hyp}.

\begin{remark}[Domain of validity]
\label{m:rem:domain}
The series \cref{m:eq:Psidef} converges for $|\zeta|<1$, and the arguments appearing
in \cref{m:eq:Gabd} are $W$ and $W\ee^{b_1t}$. In the subcritical interior regime
$\mu>\lambda$ we have $b_1<0$ and $x_-=\mu/\lambda>1$, so for $x\in[0,1]$
\begin{equation*}
0\le W\le\frac{\lambda}{\mu}<1,
\qquad
0\le W\ee^{b_1t}\le W :
\end{equation*}
the series representation is valid on the whole of the unit square, which is the
case of most interest here. When $\lambda>\mu$ the arguments may leave the unit disc
or cross the standard hypergeometric cut; one may then continue analytically,
tracking the branch along the characteristic, but the real integral
\cref{m:eq:abdconvolution} is usually safer. A convenient continuation is the Euler
integral
\begin{equation*}
{}_2F_1(1,\delta;\delta+1;\zeta) = \delta\int_0^1\frac{u^{\delta-1}}{1-\zeta u}\,\dd u
\qquad(\delta>0,\ \zeta\notin[1,\infty)),
\end{equation*}
which was used for the numerical verification quoted above. The apparent
singularity of $W$ at $x=x_-$ is a coordinate singularity only: $F(x_-,t)=x_-$ for
every $t$, and the limit of \cref{m:eq:abdconvolution} there is perfectly regular.
\end{remark}

\begin{remark}[Resonance]
\label{m:rem:resonance}
The denominators in \cref{m:eq:Psidef} vanish when $n+\delta=0$ for some integer
$n\ge1$, that is when
\begin{equation}
\label{m:eq:resonance}
\alpha = n(\mu-\lambda),\qquad n=1,2,\dots
\end{equation}
This is the classical resonance condition for linearising the $y$-equation: the
decay rate $\alpha$ of the absorption channel coincides with a harmonic of the
interior relaxation rate $\mu-\lambda$. The generating function itself is perfectly
regular there. The two $\Psi$ terms of \cref{m:eq:Gabd} each acquire a pole, but
their combination does not: in the resonant term the difference
\begin{equation*}
\frac{\bigl(W\ee^{b_1t}\bigr)^{n} - \ee^{-\alpha t}W^{n}}{n+\delta}
\end{equation*}
has a finite limit as $\delta\to-n$, producing the familiar secular factor
$t\,\ee^{-\alpha t}$ in place of a simple exponential. Equivalently
\cref{m:eq:abdconvolution} has no singular denominator at all and supplies the
limit directly. The resonant parameter values form a set of measure zero, but the
degeneracy should be kept in mind when $\alpha$ happens to sit near an integer
multiple of $\mu-\lambda$.
\end{remark}

\subsection{Recovering state probabilities}
\label{m:app:coefrecovery}

In every model above,
\begin{equation}
\label{m:eq:bivariatecoef}
p_{i,j}(t) = [x^iy^j]\,G(x,y,t)
= \frac{1}{i!\,j!}\left.\frac{\partial^{\,i+j}G}{\partial x^i\partial y^{\,j}}
\right|_{x=y=0} .
\end{equation}
For absorption--birth--death the parameter-regular representation separates the two
extractions cleanly. Writing
\begin{equation}
\label{m:eq:hdef}
h(x,t) := \alpha\ee^{-\alpha t}\int_0^t\ee^{\alpha v}F(x,v)\,\dd v,
\qquad
g(x,y,t) = \ee^{-\alpha t}y + h(x,t),
\end{equation}
the exterior variable appears only in the finite binomial expansion of $g^N$, so
\begin{equation}
\label{m:eq:abdcoef}
p_{i,j}(t) = \binom{N}{j}\ee^{-\alpha tj}\,[x^i]\,h(x,t)^{\,N-j},
\qquad 0\le j\le N .
\end{equation}
The remaining univariate coefficient can be obtained by differentiation, by
convolution of truncated series, or by Cauchy's integral formula on a circle
$|x|=\varrho<1$.

\FloatBarrier
